\chapter{Implementación}

\section{Estructura del Software}

El software del proyecto está desarrollado en Python 3.12 y utiliza ROS2 Jazzy como framework de robótica. La estructura del paquete sigue las convenciones de ROS2:

\begin{itemize}
    \item \texttt{tfg\_biped\_leg/}: Paquete principal de Python
    \begin{itemize}
        \item \texttt{odrive\_interface.py}: Clase \texttt{IOdrive} para comunicación con ODrive
        \item \texttt{odrive\_enums.py}: Definiciones de estados y errores de ODrive
        \item \texttt{odrive\_driver.py}: Wrapper de ROS2 para ODrive
        \item \texttt{leg\_kinematics.py}: Clase \texttt{LegKinematics} para cinemática
        \item \texttt{trajectory\_generator.py}: Clase \texttt{TrajectoryGenerator} para trayectorias
    \end{itemize}
    \item \texttt{scripts/}: Scripts de utilidad
    \begin{itemize}
        \item \texttt{odrive\_test.py}: Menú interactivo de pruebas
        \item \texttt{test\_single\_motor.py}: Prueba de motor único
        \item \texttt{calibrate\_motors.py}: Calibración de motores
    \end{itemize}
    \item \texttt{launch/}: Archivos de lanzamiento de ROS2
    \item \texttt{config/}: Archivos de configuración de motores
\end{itemize}

\section{Comunicación con ODrive}

La comunicación con los controladores ODrive se realiza vía USB usando la biblioteca \texttt{odrive} de Python. El proceso de inicialización incluye:

\begin{enumerate}
    \item Búsqueda y conexión con los ODrives (oficial y Makerbase)
    \item Calibración de motores y encoders
    \item Configuración de límites de corriente, velocidad y parámetros del motor
    \item Habilitación del control en lazo cerrado
\end{enumerate}

\begin{lstlisting}[language=Python, caption=Ejemplo de conexión con ODrive]
import odrive
from odrive.enums import *

# Conectar con ODrive
odrv = odrive.find_any()
print(f"ODrive conectado: {odrv.serial_number}")

# Calibrar motor en eje 0
odrv.axis0.requested_state = AXIS_STATE_FULL_CALIBRATION_SEQUENCE
while odrv.axis0.current_state != AXIS_STATE_IDLE:
    time.sleep(0.1)

# Configurar límite de corriente
odrv.axis0.motor.config.current_lim = 30  # 30A
\end{lstlisting}

\section{Cinemática Directa e Inversa}

La clase \texttt{LegKinematics} implementa las ecuaciones de cinemática para la pierna de 3-DOF. Las longitudes de eslabones son $L_1 = 0.04$m, $L_2 = 0.35$m, $L_3 = 0.35$m.

\subsection{Cinemática Directa}
Dados los ángulos articulares $\theta_1$ (Hip 1), $\theta_2$ (Hip 2), $\theta_3$ (Knee), se calcula la posición del pie:

\begin{align}
x &= L_2 \sin(\theta_2) + L_3 \sin(\theta_2 + \theta_3) \\
y &= L_1 \sin(\theta_1) + L_2 \cos(\theta_2) \sin(\theta_1) + L_3 \cos(\theta_2 + \theta_3) \sin(\theta_1) \\
z &= -[L_1 \cos(\theta_1) + L_2 \cos(\theta_2) \cos(\theta_1) + L_3 \cos(\theta_2 + \theta_3) \cos(\theta_1)]
\end{align}

\subsection{Cinemática Inversa}
Dada una posición objetivo $(x, y, z)$, se calculan los ángulos usando geometría de 3 eslabones en el espacio.

\section{Generación de Trayectorias}

La clase \texttt{TrajectoryGenerator} genera trayectorias suaves para el movimiento de la pierna, incluyendo:

\begin{itemize}
    \item Trayectorias de paso (step trajectories)
    \item Ciclos de marcha (walking cycles)
    \item Trayectorias circulares
    \item Suavizado de trayectorias usando curvas de Bézier
\end{itemize}

\section{Integración con ROS2}

El paquete de ROS2 proporciona:
\begin{itemize}
    \item Nodos para control de ODrive
    \item Topics para comandos de posición/velocidad/par
    \item Servicios para calibración y configuración
    \item Archivos de lanzamiento para iniciar el sistema
\end{itemize}

\section{Lectura de Encoders AS5600}

Los encoders magnéticos AS5600 se leen vía I2C usando el multiplexor TCA9548A. Cada encoder tiene la misma dirección I2C (0x36), por lo que el multiplexor permite seleccionar qué canal leer.

\begin{lstlisting}[language=Python, caption=Lectura de encoder AS5600 vía TCA9548A]
import smbus2
import time

class AS5600:
    def __init__(self, mux, channel, address=0x36):
        self.mux = mux
        self.channel = channel
        self.address = address
        self.bus = smbus2.SMBus(1)
    
    def read_angle(self):
        self.mux.select_channel(self.channel)
        data = self.bus.read_i2c_block_data(self.address, 0x0C, 2)
        angle = (data[0] << 8) | data[1]
        return (angle / 4096.0) * 360.0  # Grados
\end{lstlisting}
